\section*{Higham Chapter 6, Problem 6.8: Triangularization Bridge}

I am formalizing the following finite-dimensional complex matrix result in Lean
4/mathlib.

\[
\forall A\in\mathbb C^{n\times n},\ \forall \delta>0,\ \exists\hbox{ a
consistent/subordinate matrix norm } \|\cdot\| \hbox{ such that }
\|A\| \le \rho(A)+\delta .
\]

The weighted-infinity proof is already handled once the matrix is in an upper
triangular coordinate model:

\[
T\hbox{ upper triangular},\quad |T_{ii}|\le \rho
\quad\Longrightarrow\quad
\exists \nu,\ \|T\|_{\nu\to\nu}\le \rho+\delta .
\]

The remaining bridge is purely algebraic:

\[
\forall A:\mathbb C^{n\times n},\ \exists P,T,\ P\hbox{ invertible},
\quad T=P^{-1}AP,\quad T\hbox{ upper triangular},\quad
\forall i,\ |T_{ii}|\le \rho(A).
\]

In the local Lean development, the conditional theorem consumes:

\[
T = P^{-1} A P,\qquad T\hbox{ is upper triangular},\qquad
\forall i,\ |T_{ii}|\le \rho .
\]

Please give a Lean-friendly route for proving this bridge in current mathlib.
Important available-looking facts/APIs from mathlib search:

\begin{itemize}
\item \texttt{Module.End.exists\_eigenvalue} in
\texttt{Mathlib.LinearAlgebra.Eigenspace.Triangularizable}: every nontrivial
finite-dimensional endomorphism over an algebraically closed field has an
eigenvalue.
\item \texttt{Module.End.iSup\_maxGenEigenspace\_eq\_top}: generalized
eigenspaces span the whole space over an algebraically closed field.
\item \texttt{Matrix.charpoly\_of\_upperTriangular}: an upper triangular matrix
has characteristic polynomial \(\prod_i (X-T_{ii})\).
\item \texttt{Module.End.HasEigenvalue.mem\_spectrum} and
\texttt{hasEigenvalue\_iff\_mem\_spectrum}.
\item \texttt{Matrix.BlockTriangular id} is the apparent upper-triangular
predicate used by mathlib.
\item I did not find a ready-made theorem named ``triangularizable gives an
upper triangular basis/matrix''.
\end{itemize}

Question:

\begin{enumerate}
\item Is there a ready-made mathlib declaration that directly gives a basis in
which a complex finite-dimensional endomorphism is upper triangular? If yes,
please name the declarations and the proof outline.
\item If not, what is the smallest Lean-friendly bridge to formalize? I prefer
a theorem stated in terms of an existence of a basis \(b\) such that
\(\operatorname{LinearMap.toMatrix}\ b\ b\ f\) is
\texttt{Matrix.BlockTriangular id}. Give the induction structure and the key
submodule/basis-extension lemmas to search for.
\item Once such a triangular basis exists, how should I prove that each diagonal
entry is an eigenvalue or at least a root of the characteristic polynomial, and
therefore bounded by an abstract spectral-radius maximum predicate?
\item Please flag any hidden assumptions: nonempty dimensions, zero dimension,
choice of \(\rho(A)\) as a supplied maximum over eigenvalue moduli, and whether
the statement should use eigenvalues, spectrum, or charpoly roots to connect
diagonal entries to \(\rho(A)\).
\end{enumerate}

Please answer as a concrete Lean implementation plan, not prose only. Avoid
depending on unproved Schur decomposition or unitary triangularization: ordinary
algebraic triangularization is enough.
